\SourcePage{533}{536}
\section{$6j$ symbols}

In \S~106 we defined the $3j$ symbols as the coefficients in the sum (106.4), which represents the wavefunction of a three-particle system with zero total angular momentum. From the standpoint of transformation properties under rotations, this sum is a scalar. It follows that the set of $3j$ symbols for specified values of $j_1,j_2,j_3$ (and all possible $m_1,m_2,m_3$) may be regarded as a collection of quantities which, under rotations, transform contragrediently to the products $\psi_{j_1m_1}\psi_{j_2m_2}\psi_{j_3m_3}$, thereby making the entire sum invariant.

From this viewpoint, one may ask how to construct a scalar solely from $3j$-symbols. Such a scalar may depend only on the $j$ values, and not on the

\SourcePage{534}{537}
$m$ values, which change under rotations. In other words, it must be expressible through sums over all the $m$ values. Each such summation amounts to a ``contraction'' of the product of two $3j$-symbols according to the rule

\begin{equation}
 \sum_m(-1)^{j-m}
 \begin{pmatrix}j&\cdots\\m&\cdots\end{pmatrix}
 \begin{pmatrix}j&\cdots\\-m&\cdots\end{pmatrix}
 \tag{108.1}
\end{equation}
(compare the construction of the scalar (106.2)).

Since every ``contraction'' involves a pair of numbers $m$, a complete scalar must be constructed from a product of an even number of $3j$ symbols. By their orthogonality property, contracting a product of two $3j$ symbols gives the trivial result
\[
\begin{aligned}
&\sum_{m_1m_2m_3}
\begin{pmatrix}j_1&j_2&j_3\\m_1&m_2&m_3\end{pmatrix}
\begin{pmatrix}j_1&j_2&j_3\\-m_1&-m_2&-m_3\end{pmatrix}
(-1)^{j_1+j_2+j_3-m_1-m_2-m_3}\\
&\qquad=\sum_{m_1m_2m_3}
\begin{pmatrix}j_1&j_2&j_3\\m_1&m_2&m_3\end{pmatrix}^{2}=1
\end{aligned}
\]
(here the equality $m_1+m_2+m_3=0$ and formulas (106.6) and (106.12) have been used). Hence the smallest number of factors required to obtain a nontrivial scalar is four. In each $3j$ symbol, the three numbers $j$ geometrically form a closed triangle. Since every number $j$ must occur in two $3j$ symbols during the contraction, it is clear that a scalar formed from a product of four $3j$ symbols contains six numbers $j$. Geometrically, they must be represented by the edge lengths of an irregular tetrahedron (Fig.~45), with one face corresponding to each $3j$ symbol. A definite convention for carrying out the contractions is adopted in defining the scalar in question; it is expressed by the formula
\begin{equation}
\begin{aligned}
\left\{\begin{matrix}j_1&j_2&j_3\\j_4&j_5&j_6\end{matrix}\right\}
={}&\sum_{\text{all }m}(-1)^{\sum_i(j_i-m_i)}
\begin{pmatrix}j_1&j_2&j_3\\-m_1&-m_2&-m_3\end{pmatrix}
\begin{pmatrix}j_1&j_5&j_6\\m_1&-m_5&m_6\end{pmatrix}\\
&\times\begin{pmatrix}j_4&j_2&j_6\\m_4&m_2&-m_6\end{pmatrix}
\begin{pmatrix}j_4&j_5&j_3\\-m_4&m_5&m_3\end{pmatrix}.
\end{aligned}
\tag{108.2}
\end{equation}
The summation here is over all possible values of every number $m$; however, since the sum of the three $m$'s in each $3j$ symbol must vanish, only

\SourcePage{535}{538}
three of the six $m$'s are actually independent. The quantities defined by (108.2) are called $6j$ symbols or Racah coefficients\footnote{The notation
\[
W(j_1j_2j_4j_5;j_3j_6)=(-1)^{j_1+j_2+j_4+j_5}
\left\{\begin{matrix}j_1&j_2&j_3\\j_4&j_5&j_6\end{matrix}\right\}
\]
is also used in the literature.}.

From definition (108.2), with account taken of the symmetry properties of the $3j$ symbols, it is readily verified that the $6j$ symbol is unchanged under any permutation of its three columns, and that the upper and lower numbers may be interchanged in any pair of columns. By virtue of these symmetry properties, the sequence of numbers $j_1\ldots j_6$ in a $6j$ symbol can be represented in 24 equivalent forms\footnote{If the tetrahedron in Fig.~45 is imagined to be regular, the 24 equivalent permutations of the numbers $j$ can be obtained by applying the 24 symmetry transformations (rotations and reflections) of the tetrahedron.}.


\begin{figure}[htbp]
\centering
\begin{tikzpicture}[x=1cm,y=1cm,line width=.45pt,line cap=round,font=\small]
 \coordinate (L) at (-1.55,0.45); \coordinate (R) at (1.55,0.45);
 \coordinate (T) at (0.55,1.65); \coordinate (B) at (-0.65,-1.85);
 \draw (L)--(R)--(T)--(L)--(B)--(R);
 \draw[dashed] (T)--(B);
 \node[above left] at ($(L)!0.56!(T)$) {$j_1$};
 \node[above right] at ($(T)!0.48!(R)$) {$j_2$};
 \node[below] at ($(L)!0.40!(R)$) {$j_3$};
 \node[right] at ($(R)!0.48!(B)$) {$j_4$};
 \node[left] at ($(L)!0.50!(B)$) {$j_5$};
 \node[right] at ($(T)!0.55!(B)$) {$j_6$};
\end{tikzpicture}
\caption{}
\label{fig:45}
\end{figure}

The $6j$ symbols also have another, less evident symmetry property, which equates symbols having different sets of numbers $j$:
\begin{equation}
\left\{\begin{matrix}j_1&j_2&j_3\\j_4&j_5&j_6\end{matrix}\right\}=
\left\{\begin{matrix}
j_1&\frac12(j_2+j_5+j_3-j_6)&\frac12(j_3+j_6+j_2-j_5)\\
j_4&\frac12(j_2+j_5+j_6-j_3)&\frac12(j_3+j_6+j_5-j_2)
\end{matrix}\right\}.
\tag{108.3}
\end{equation}
($T.\ Regge$, 1959)\footnote{See the note on p.~529.}.

We give a useful relation between $6j$ and $3j$ symbols that follows from definition (108.2):
\begin{equation}
\begin{aligned}
&\sum_{m_4m_5m_6}(-1)^{j_4+j_5+j_6-m_4-m_5-m_6}
\begin{pmatrix}j_1&j_5&j_6\\m_1&-m_5&m_6\end{pmatrix}
\begin{pmatrix}j_4&j_2&j_6\\m_4&m_2&-m_6\end{pmatrix}\\
&\qquad\times\begin{pmatrix}j_4&j_5&j_3\\-m_4&m_5&m_3\end{pmatrix}
=\begin{pmatrix}j_1&j_2&j_3\\m_1&m_2&m_3\end{pmatrix}
\left\{\begin{matrix}j_1&j_2&j_3\\j_4&j_5&j_6\end{matrix}\right\}.
\end{aligned}
\tag{108.4}
\end{equation}
The expression being summed on the left-hand side differs from the sum in (108.2) by the omission of one factor (a $3j$-symbol). Thus the sum in (108.4) may be represented by a tetrahedron (Fig.~45) with one face removed; this is why the sum is not a scalar. In other words, in its transfor\SourcePageJoin{536}{539}mation properties it corresponds to a single $3j$-symbol---the one on the right-hand side of (108.4)---and must therefore be proportional to that symbol. The proportionality factor (the $6j$-symbol on the right-hand side) is obtained at once by multiplying both sides by
\[
\begin{pmatrix}j_1&j_2&j_3\\m_1&m_2&m_3\end{pmatrix}
\]
and summing over the remaining numbers $m_1,m_2,m_3$.


$6j$-symbols arise naturally in the following question concerning the addition of three angular momenta.

Let the three angular momenta $j_1,j_2,j_3$ combine to form a resultant angular momentum $J$. Specifying $J$ (and its projection $M$), however, does not yet determine the state of the system uniquely: the state also depends on the order in which the angular momenta are added (or, as it is called, their coupling scheme).


Consider, for example, the following two coupling schemes: 1) first $j_1$ and $j_2$ are added to form the total angular momentum $j_{12}$, after which $j_{12}$ and $j_3$ are added to give the final angular momentum $J$; 2) $j_2$ and $j_3$ are added to give $j_{23}$, and then $j_{23}$ is combined with $j_1$ to give $J$. The first scheme corresponds to states in which, besides $j_1,j_2,j_3,J,M$, the quantity $j_{12}$ has a definite value; let their wave functions be denoted by $\psi_{j_{12}JM}$ (the repeated indices $j_1j_2j_3$ being suppressed for brevity). Similarly, denote the wave functions for the second coupling scheme by $\psi_{j_{23}JM}$. In either case the ``intermediate'' angular momentum ($j_{12}$ or $j_{23}$) can, in general, take more than one value. Thus, for fixed $J,M$, there are two distinct sets of states, distinguished respectively by the values of $j_{12}$ and $j_{23}$. By the general rules, the functions in these two sets are related by a certain unitary transformation

\begin{equation}
 \psi_{j_{23}JM}=\sum_{j_{12}}\langle j_{12}|j_{23}\rangle\psi_{j_{12}JM}.
 \tag{108.5}
\end{equation}
It is clear on physical grounds that the coefficients of this transformation do not depend on $M$: they cannot depend on the orientation of the entire system in space. They therefore depend only on the six angular momenta $j_1,j_2,j_3,j_{12},j_{23},J$, not on their projections, and so are scalar quantities (in the sense specified above). These coefficients can be calculated directly as follows.

\SourcePage{537}{540}
Applying formula (106.9) twice, we find
\[
\begin{aligned}
\psi_{j_{23}JM}&=\sum_{(m)}\langle m_1m_{23}|JM\rangle\psi_{j_1m_1}\psi_{j_{23}m_{23}}\\
&=\sum_{(m)}\langle m_1m_{23}|JM\rangle\langle m_2m_3|j_{23}m_{23}\rangle
\psi_{j_1m_1}\psi_{j_2m_2}\psi_{j_3m_3},\\
\psi_{j_{12}JM}&=\sum_{(m)}\langle m_3m_{12}|JM\rangle\langle m_1m_2|j_{12}m_{12}\rangle
\psi_{j_1m_1}\psi_{j_2m_2}\psi_{j_3m_3}
\end{aligned}
\]
(the sign $(m)$ under the summation sign means that the summation is over all the numbers $m_1,m_2,\ldots$ occurring in the expression). Using the orthonormality of the functions $\psi_{jm}$, we now obtain
\[
\begin{aligned}
\langle j_{12}|j_{23}\rangle&\equiv\int\psi_{j_{12}JM}^{*}\psi_{j_{23}JM}\,dq\\
&=\sum_{(m)}\langle m_3m_{12}|JM\rangle\langle m_1m_{23}|JM\rangle
\langle m_1m_2|j_{12}m_{12}\rangle\langle m_2m_3|j_{23}m_{23}\rangle.
\end{aligned}
\]
The sum on the right is taken at a specified value of $M$, but the result of the summation is in fact independent of $M$, for the reason given above. The summation may therefore be extended over the values of $M$ as well, provided that the factor $1/(2J+1)$ is placed before the sum. Then, expressing the coefficients $\langle m_1m_2|jm\rangle$ in terms of $3j$ symbols according to (106.10), we obtain
\begin{equation}
\langle j_{12}|j_{23}\rangle=(-1)^{j_1+j_2+j_3+J}\sqrt{(2j_{12}+1)(2j_{23}+1)}
\left\{\begin{matrix}j_1&j_2&j_{12}\\j_3&J&j_{23}\end{matrix}\right\}.
\tag{108.6}
\end{equation}
The connection of the $6j$ symbols with the transformation coefficients (108.5) makes it easy to derive several useful formulas for sums of products of $6j$ symbols.

First, by the unitarity of transformation (108.5) (and the reality of its coefficients), we have
\begin{equation}
\sum_j(2j+1)(2j''+1)
\left\{\begin{matrix}j_1&j_2&j'\\j_3&j_4&j\end{matrix}\right\}
\left\{\begin{matrix}j_3&j_2&j\\j_1&j_4&j''\end{matrix}\right\}
=\delta_{j'j''}.
\tag{108.7}
\end{equation}
Next consider three schemes for coupling three angular momenta, whose intermediate sums are respectively $j_{12},j_{23}$, and $j_{31}$. By the rule for matrix multiplication, the corresponding transformation coefficients (108.6) are related by
\[
\sum_{j_{23}}\langle j_{12}|j_{23}\rangle\langle j_{23}|j_{31}\rangle=\langle j_{12}|j_{31}\rangle.
\]

\SourcePage{538}{541}
Substitution of (108.6), followed by a change of index notation, gives
\begin{equation}
\sum_j(-1)^{j+j_3+j_6}(2j+1)
\left\{\begin{matrix}j_2&j_4&j_6\\j_1&j_5&j\end{matrix}\right\}
\left\{\begin{matrix}j_4&j_1&j_6\\j_2&j_5&j_3\end{matrix}\right\}
=\left\{\begin{matrix}j_1&j_2&j_3\\j_4&j_5&j_6\end{matrix}\right\}.
\tag{108.8}
\end{equation}
Finally, consideration of different schemes for coupling four angular momenta yields\footnote{See the book by Edmonds cited above.} the following addition formula for products of three $6j$ symbols:
\begin{equation}
\begin{aligned}
&\sum_j(-1)^{j+\sum_{1}^{9}j_i}(2j+1)
\left\{\begin{matrix}j_4&j_2&j_6\\j_9&j_8&j\end{matrix}\right\}
\left\{\begin{matrix}j_2&j_1&j_3\\j_7&j&j_9\end{matrix}\right\}
\left\{\begin{matrix}j_4&j_3&j_5\\j_7&j_8&j\end{matrix}\right\}\\
&\qquad=\left\{\begin{matrix}j_1&j_2&j_3\\j_4&j_5&j_6\end{matrix}\right\}
\left\{\begin{matrix}j_6&j_1&j_5\\j_7&j_8&j_9\end{matrix}\right\}.
\end{aligned}
\tag{108.9}
\end{equation}
($L.\ C.\ Biedenharn$, $J.\ P.\ Elliott$, 1953).

For reference, we give some explicit expressions for $6j$ symbols. In the general case a $6j$ symbol can be represented by the following sum:

\begingroup\footnotesize
\begin{equation}
\begin{aligned}
\left\{\begin{matrix}j_1&j_2&j_3\\j_4&j_5&j_6\end{matrix}\right\}
={}&\Delta(j_1j_2j_3)\Delta(j_1j_5j_6)\Delta(j_4j_2j_6)\Delta(j_4j_5j_3)\\
&\times\sum_z\frac{(-1)^z(z+1)!}{(z-j_1-j_2-j_3)!(z-j_1-j_5-j_6)!(z-j_4-j_2-j_6)!(z-j_4-j_5-j_3)!}\\
&\times(j_1+j_2+j_4+j_5-z)!(j_2+j_3+j_5+j_6-z)!(j_3+j_1+j_6+j_4-z)! ,
\end{aligned}
\tag{108.10}
\end{equation}
\endgroup
where
\[
\Delta(abc)=\left[\frac{(a+b-c)!(a-b+c)!(-a+b+c)!}{(a+b+c+1)!}\right]^{1/2},
\]
The sum runs over every positive integral value of $z$ for which no factorial in the denominator has a negative argument. Table~10 gives formulas for the $6j$-symbols when one parameter is 0, $1/2$, or 1.


We conclude with several remarks about higher-order scalars formed from $3j$ symbols.

After the $6j$ symbol, the next in complexity is the scalar formed by contracting a product of six $3j$ symbols. These $3j$ symbols contain 18 pairwise equal numbers $j$,

\SourcePage{539}{542}
so that the resulting scalar depends on nine parameters $j$. It is conventionally called a $9j$ symbol and is defined as follows ($E.\ Wigner$, 1951)\footnote{According to the general contraction rule (108.1), the arguments $m$ in the last three $3j$ symbols should be written with minus signs, and the factor $(-1)^{\sum(j-m)}$ should be inserted under the summation sign. Using property (106.6) of the $3j$ symbols, however, and noting that in this case, as is readily seen, the sum $\sum m$ over all nine numbers $m$ is zero, we arrive at definition (108.11).}:
\begin{equation}
\begin{aligned}
\left\{\begin{matrix}j_{11}&j_{12}&j_{13}\\j_{21}&j_{22}&j_{23}\\j_{31}&j_{32}&j_{33}\end{matrix}\right\}
={}&\sum_{\text{all }m}
\begin{pmatrix}j_{11}&j_{12}&j_{13}\\m_{11}&m_{12}&m_{13}\end{pmatrix}
\begin{pmatrix}j_{21}&j_{22}&j_{23}\\m_{21}&m_{22}&m_{23}\end{pmatrix}\\
&\times\begin{pmatrix}j_{31}&j_{32}&j_{33}\\m_{31}&m_{32}&m_{33}\end{pmatrix}
\begin{pmatrix}j_{11}&j_{21}&j_{31}\\m_{11}&m_{21}&m_{31}\end{pmatrix}\\
&\times\begin{pmatrix}j_{12}&j_{22}&j_{32}\\m_{12}&m_{22}&m_{32}\end{pmatrix}
\begin{pmatrix}j_{13}&j_{23}&j_{33}\\m_{13}&m_{23}&m_{33}\end{pmatrix}.
\end{aligned}
\tag{108.11}
\end{equation}

\begin{table}[p]
\centering\small
\caption{Formulas for $6j$ symbols}
\label{tab:10}
\renewcommand{\arraystretch}{1.55}
\begin{tabular}{@{}l@{}}
\hline
$\displaystyle \left\{\begin{matrix}a&b&c\\0&c&b\end{matrix}\right\}=\frac{(-1)^s}{\sqrt{(2b+1)(2c+1)}},\qquad s=a+b+c$\\ \hline
$\displaystyle \left\{\begin{matrix}a&b&c\\\frac12&c-\frac12&b+\frac12\end{matrix}\right\}=(-1)^s\left[\frac{(s-2b)(s-2c+1)}{(2b+1)(2b+2)2c(2c+1)}\right]^{1/2}$\\
$\displaystyle \left\{\begin{matrix}a&b&c\\\frac12&c-\frac12&b-\frac12\end{matrix}\right\}=(-1)^s\left[\frac{(s+1)(s-2a)}{2b(2b+1)2c(2c+1)}\right]^{1/2}$\\ \hline
$\displaystyle \left\{\begin{matrix}a&b&c\\1&c-1&b-1\end{matrix}\right\}=(-1)^s\left[\frac{s(s+1)(s-2a-1)(s-2a)}{(2b-1)2b(2b+1)(2c-1)2c(2c+1)}\right]^{1/2}$\\
$\displaystyle \left\{\begin{matrix}a&b&c\\1&c-1&b\end{matrix}\right\}=(-1)^s\left[\frac{2(s+1)(s-2a)(s-2b)(s-2c+1)}{2b(2b+1)(2b+2)(2c-1)2c(2c+1)}\right]^{1/2}$\\
$\displaystyle \left\{\begin{matrix}a&b&c\\1&c-1&b+1\end{matrix}\right\}=(-1)^s\left[\frac{(s-2b-1)(s-2b)(s-2c+1)(s-2c+2)}{(2b+1)(2b+2)(2b+3)(2c-1)2c(2c+1)}\right]^{1/2}$\\
$\displaystyle \left\{\begin{matrix}a&b&c\\1&c&b\end{matrix}\right\}=(-1)^{s+1}\frac{2[b(b+1)+c(c+1)-a(a+1)]}{[2b(2b+1)(2b+2)2c(2c+1)(2c+2)]^{1/2}}$\\ \hline
\end{tabular}
\end{table}

This quantity can also be represented as a sum of

\SourcePage{540}{543}
products of three $6j$ symbols:
\begin{equation}
\begin{aligned}
\left\{\begin{matrix}j_{11}&j_{12}&j_{13}\\j_{21}&j_{22}&j_{23}\\j_{31}&j_{32}&j_{33}\end{matrix}\right\}
={}&\sum_j(-1)^{2j}(2j+1)
\left\{\begin{matrix}j_{11}&j_{21}&j_{31}\\j_{32}&j_{33}&j\end{matrix}\right\}\\
&\times\left\{\begin{matrix}j_{12}&j_{22}&j_{32}\\j_{21}&j&j_{23}\end{matrix}\right\}
\left\{\begin{matrix}j_{13}&j_{23}&j_{33}\\j&j_{11}&j_{12}\end{matrix}\right\}.
\end{aligned}
\tag{108.12}
\end{equation}
The equivalence of (108.11) and (108.12) may be checked by inserting definition (108.2) into (108.12) and using the orthogonality properties of the $3j$-symbols.

The high symmetry of the $9j$-symbol follows directly from definition (108.11) and the symmetry properties of the $3j$-symbols. Interchanging any two rows or any two columns multiplies the $9j$-symbol by $(-1)^{\sum j}$. It is also unchanged by transposition, that is, by interchanging rows and columns.

Scalars of still higher order depend on still more parameters $j$. Their number must clearly be a multiple of three; these are the $3nj$-symbols. We shall not consider their properties here. We mention only that for every $n>3$ there is more than one distinct type of $3nj$-symbol, and these types cannot be reduced to one another. Thus there are two distinct types of $12j$-symbols\footnote{For a fuller account of the theory of $9j$-symbols and of the properties of $3nj$-symbols, see the book by Edmonds cited on p.~272, and the books: A.~P.~Yutsis, I.~B.~Levinson, V.~V.~Vanagas. \emph{Mathematical Apparatus of the Theory of Angular Momentum}. --- Vilnius, 1960; D.~A.~Varshalovich, A.~N.~Moskalev, V.~K.~Khersonskii. \emph{Quantum Theory of Angular Momentum}. --- Moscow: Nauka, 1975.}.

